
\documentclass[12pt]{exam}
\input{../tmpl/packages.tex}

% Theorem environments
\newtheorem*{exercise}{Exercise}
\newenvironment{exe}[1]{\begin{exercise}[#1\label{#1}]}{\end{exercise}}
\newcommand{\mtexe}[1]{\noindent\textbf{Exercise #1.}}

% Script letter shorthands
\newcommand{\scA}{\mathscr{A}}
\newcommand{\scF}{\mathscr{F}}
\newcommand{\scG}{\mathscr{G}}
\newcommand{\scH}{\mathscr{H}}
\newcommand{\scI}{\mathscr{I}}
\newcommand{\scJ}{\mathscr{J}}
\newcommand{\scL}{\mathscr{L}}
\newcommand{\scM}{\mathscr{M}}
\newcommand{\scO}{\mathscr{O}}
\newcommand{\scZ}{\mathscr{Z}}

% Ideals and other fraktures
\newcommand{\frd}{\mathfrak{d}}
\newcommand{\frm}{\mathfrak{m}}
\newcommand{\frp}{\mathfrak{p}}

% Blackboard style
\renewcommand{\AA}{\mathbb{A}}
\newcommand{\CC}{\mathbb{C}}
\newcommand{\FF}{\mathbb{F}}
\newcommand{\NN}{\mathbb{N}}
\newcommand{\PP}{\mathbb{P}}
\newcommand{\ZZ}{\mathbb{Z}}

% Functions
\newcommand{\Cl}{\operatorname{Cl}}
\newcommand{\Frac}{\operatorname{Frac}}
\newcommand{\Hom}{\operatorname{Hom}}
\newcommand{\Pic}{\operatorname{Pic}}
\newcommand{\Proj}{\operatorname{Proj}}
\newcommand{\Spec}{\operatorname{Spec}}
\newcommand{\id}{\operatorname{id}}
\newcommand{\im}{\operatorname{im}}
\newcommand{\scHom}{\operatorname{\mathscr{H}om}}

% Other
\newcommand*{\twomat}[4]{\left(\begin{array}{cc} #1 & #2 \\ #3 & #4 \end{array}\right)}
\newcommand*{\prob}[1]{\hyperref[#1]{(#1)}}



\begin{document}
\pagestyle{plain}
\thispagestyle{empty}

\noindent
\begin{tabular*}{\textwidth}{l @{\extracolsep{\fill}} r @{\extracolsep{6pt}} l}
\textbf{Vakil - The Rising Sea} & \textbf{Name:} & Abhay Goel \\
\textbf{Solutions} & \textbf{Last updated:} & \today \\
\end{tabular*}\\
\rule[2ex]{\textwidth}{2pt}

\section*{Chapter 1}

\mtexe{1.2.A}
\begin{proof}
	The two notions can be identified. Indeed, given a group $G$ in the usual sense, define a category with a single object $X$, and for each $g \in G$, define a morphism $f_g : X \to X$. Further, for two morphisms $f_g,f_h$, define composition by $f_g \circ f_h = f_{gh}$. This gives a 1-object groupoid, since each morphism is an isomorphism: the morphism $f_g$ has inverse $f_{g^{-1}}$.
	
	Conversely, given a 1-object groupoid, the set of morphisms on that one object forms a group. \\
	
	On the other hand, there are clearly groupoids that aren't groups, e.g. a category with two objects and only the identity morphisms on those objects.
\end{proof}

\mtexe{1.2.B}
\begin{proof}
	The identity morphism is invertible, with itself as left- and right-inverse. The composition of two invertible morphisms is again invertible, since $(f \circ g) \circ (g^{-1} \circ f^{-1}) = \id$. Finally, the inverse of an invertible morphism is invertible. So, the set of invertible morphisms forms a group under composition. \\
	
	Given a set $X$, the automorphisms of $X$ in the category of sets is the set of bijections $f : X \to X$, i.e. the automorphism group is $S_X$, the symmetric group on $X$. The automorphisms of a vector space in $Vec_k$ are the automorphisms in the usual sense in linear algebra. \\
	
	Finally, if $A$ and $B$ are isomorphic objects, with an isomorphism $f : A \to B$, then for an automorphism $\phi : A \to A$, the morphism $f \circ \phi \circ f^{-1}$ is an automorphism of $B$, with inverse $f \circ \phi^{-1} \circ f^{-1}$. So, this association gives an isomorphism of automorphism groups.
\end{proof}

\mtexe{1.2.C} Skipped upon recommendation

\mtexe{1.2.D} Skipped upon recommendation

\mtexe{1.3.A}
\begin{proof}
	Let $A,B$ both be initial objects in some category. Then, since $A$ is initial, there is a morphism $f : A \to B$, and since $B$ is initial, there is a morphism $g : B \to A$. Then $\id_A$ and $g \circ f$ are both morphisms $A \to A$, but since $A$ is initial, there is a unique such morphism, so $g \circ f = \id_A$. Similarly, $f \circ g = \id_B$. So, $f$ and $g$ are isomorphisms. Further, this isomorphism is unique since any isomorphism $A \to B$ is firstly a morphism, and there is a unique such morphism since $A$ is initial.
	
	The case of final objects is identical. In fact, this proof already shows it by noting that an object is final in a category if and only if it is initial in the opposite category.
\end{proof}

\mtexe{1.3.B}
\begin{proof}[Solution]
	The initial object in $Sets$ is the empty set, since there is a unique map from the empty set to any other set. The final object in $Sets$ is any set with a single element, since there is a unique map from any set to this set. Note that specifying which single-element set is irrelevant, since any two such sets are uniquely isomorphic as sets. \\
	
	The initial object in $Rings$ is $\ZZ$ (recall we are only considering commutative rings with unity and morphisms that preserve unities), since there is a unique morphism given by mapping $1$ to $1 \in R$ for a given ring $R$. The final object is the zero ring. \\
	
	$Top$ is the ``same'' as set: the initial object is the empty space and the final object is a single-point space, since there is a unique continuous map in each case. \\
	
	Given a set $S$, the poset category of subsets of $S$ has $\emptyset$ as an initial object since it is a subset of each other subset of $S$, and it has $S$ as a final object. If $S$ is a topological space and we consider the poset category of open subsets of $S$, the initial and final objects are the same.
\end{proof}

\mtexe{1.3.C}
\begin{proof}
	In one direction, suppose $s \in S$ is a zerodivisor, so that there is some nonzero $x \in A$ with $sx = 0$. Then, since $s(1 \cdot x - 1 \cdot 0)$, we get that $x/1 = 0/1$ in $S^{-1}A$. So, $x$ is in the kernel of the map $A \to S^{-1}A$ and it isn't injective.
	
	Conversely, suppose this map isn't injective. Then, there is some nonzero element $x$ in the kernel, so that $x/1 = 0/1$. I.e. there is some $s \in S$ with $s(1 \cdot x - 1 \cdot 0) = 0$, i.e. $sx = 0$, whence $s$ is a zerodivisor.
\end{proof}

\mtexe{1.3.D}
\begin{proof}
	First, note that $A \to S^{-1}A$ has the property that all elements of $S$ are mapped to invertible elements, since for $s \in S$, the image is $s/1$, which has inverse $1/s \in S^{-1}A$. Now we want to see that it is initial with respect to this property. So, let $f : A \to B$ be a morphism with each element of $S$ mapping to an invertible element of $B$. Then, we'd like to define a factorization through $S^{-1}A$. Define:
	\[ g(a/s) = f(a)f(s)^{-1} \]
	First, note that this definition makes sense, since for each $s \in S$, $f(s)$ is invertible and so $f(s)^{-1}$ exists. Second, we want to see that it is well-defined. So, suppose $a/s = b/t$, so that there is some $r \in S$ with $r(ta-sb) = 0$. Then,
	\[ 0 = f(r(ta-sb)) = f(r)(f(t)f(a)-f(s)f(b)) \]
	Since $f(r)$ is invertible, we can cancel it, and then multiplying by the inverses of $f(t)$ and $f(s)$ gives:
	\[ f(a)f(s)^{-1} = f(b)f(t)^{-1} \]
	so that $g$ is well-defined. Finally, we want to see it is a ring homomorphism. First, we have $g(1/1) = f(1)f(1)^{-1} = 1$. Second, if $a/s,b/t \in S^{-1}A$, then
	\[ g(a/s+b/t) = g((at+bs)/(st)) = f(at+bs)f(st)^{-1} = f(a)f(s)^{-1}+f(b)f(t)^{-1} = g(a/s)+g(b/t) \]
	and
	\[ g((a/s)(b/t)) = g((ab)/(st)) = f(ab)f(st)^{-1} = f(a)f(b)f(s)^{-1}f(t)^{-1} = g(a/s)g(b/t) \]
	as desired. Thus $g$ is a ring homomorphism. Penultimately, we show it factors $f$. Let $\iota : A \to S^{-1}A$ be the canonical map. Then, for $a \in A$,
	\[ g(\iota(a)) = g(a/1) = f(a)f(1)^{-1} = f(a) \]
	so that $f = g \circ \iota$ as desired. To finally finish, we need to see that $g$ is unique. So, suppose we have another ring homomorphism $h : S^{-1}A \to B$ with $f = h \circ \iota$. Then, let $a/s \in S^{-1}A$. We have
	\[ h(a/s) = h((a/1)(1/s)) = h(a/1)h(1/s) = h(a/1)h(s/1)^{-1} = h(\iota(a))h(\iota(s))^{-1} = f(a)f(s)^{-1} = g(a/s) \]
	so that $h = g$. This completes the proof that $S^{-1}A$ is initial among such rings. \\
	
	We address the ``furthermore.'' Let $M$ be an $S^{-1}A$-module. Then, it is also an $A$-module via $a \cdot m = (a/1) \cdot m$ for $a \in A$ and $m \in M$ (restriction of scalars). For $s \in S$, multiplication by $s$ is an isomorphism of $M$ since it has inverse given by multiplication by $1/s$.
	
	Conversely, suppose $M$ is an $A$ module such that multiplication by $s$ is an automorphism of $M$ for all $s \in S$. For $x \in A$, let $\varphi(x) : M \to M$ denote multiplication by $x$, and let $R' = \varphi(A)$ be the image of $A$ in the endomorphism ring of $M$. Then $R'$ is a commutative ring with unity, and an application of Zorn's Lemma gives a maximal subring $R$ of the endomorphism ring of $M$ containing $R'$. Then, $\varphi : A \to R$ and each element of $S$ maps to an invertible element of $R$, so this factors as a map from $S^{-1}A$, i.e. we get a map from $S^{-1}A$ into the endomorphism ring of $M$, i.e. $M$ is an $S^{-1}A$-module.
\end{proof}

\end{document}